\documentclass{amsart}
% For dealing with references we use the comment environment
\usepackage{verbatim}
\newenvironment{reference}{\comment}{\endcomment}
%\newenvironment{reference}{}{}
\newenvironment{slogan}{\comment}{\endcomment}
\newenvironment{history}{\comment}{\endcomment}

% For commutative diagrams we use Xy-pic
\usepackage[all]{xy}

% We use 2cell for 2-commutative diagrams.
\xyoption{2cell}
\UseAllTwocells

% We use multicol for the list of chapters between chapters
\usepackage{multicol}

% This is generally recommended for better output
\usepackage{lmodern}
\usepackage[T1]{fontenc}

% For cross-file-references

% Package for hypertext links:
\usepackage{hyperref}

% For any local file, say "hello.tex" you want to link to please

% Theorem environments.
%
\theoremstyle{plain}
\newtheorem{theorem}[subsection]{Theorem}
\newtheorem{proposition}[subsection]{Proposition}
\newtheorem{lemma}[subsection]{Lemma}

\theoremstyle{definition}
\newtheorem{definition}[subsection]{Definition}
\newtheorem{example}[subsection]{Example}
\newtheorem{exercise}[subsection]{Exercise}
\newtheorem{situation}[subsection]{Situation}

\theoremstyle{remark}
\newtheorem{remark}[subsection]{Remark}
\newtheorem{remarks}[subsection]{Remarks}

\numberwithin{equation}{subsection}

% Macros
%
\def\lim{\mathop{\mathrm{lim}}\nolimits}
\def\colim{\mathop{\mathrm{colim}}\nolimits}
\def\Spec{\mathop{\mathrm{Spec}}}
\def\Hom{\mathop{\mathrm{Hom}}\nolimits}
\def\Ext{\mathop{\mathrm{Ext}}\nolimits}
\def\SheafHom{\mathop{\mathcal{H}\!\mathit{om}}\nolimits}
\def\SheafExt{\mathop{\mathcal{E}\!\mathit{xt}}\nolimits}
\def\Sch{\mathit{Sch}}
\def\Mor{\mathop{\mathrm{Mor}}\nolimits}
\def\Ob{\mathop{\mathrm{Ob}}\nolimits}
\def\Sh{\mathop{\mathit{Sh}}\nolimits}
\def\NL{\mathop{N\!L}\nolimits}
\def\CH{\mathop{\mathrm{CH}}\nolimits}
\def\proetale{{pro\text{-}\acute{e}tale}}
\def\etale{{\acute{e}tale}}
\def\QCoh{\mathit{QCoh}}
\def\Ker{\mathop{\mathrm{Ker}}}
\def\Im{\mathop{\mathrm{Im}}}
\def\Coker{\mathop{\mathrm{Coker}}}
\def\Coim{\mathop{\mathrm{Coim}}}

% Boxtimes
%
\DeclareMathSymbol{\boxtimes}{\mathbin}{AMSa}{"02}

%
% Macros for moduli stacks/spaces
%
\def\QCohstack{\mathcal{QC}\!\mathit{oh}}
\def\Cohstack{\mathcal{C}\!\mathit{oh}}
\def\Spacesstack{\mathcal{S}\!\mathit{paces}}
\def\Quotfunctor{\mathrm{Quot}}
\def\Hilbfunctor{\mathrm{Hilb}}
\def\Curvesstack{\mathcal{C}\!\mathit{urves}}
\def\Polarizedstack{\mathcal{P}\!\mathit{olarized}}
\def\Complexesstack{\mathcal{C}\!\mathit{omplexes}}
% \Pic is the operator that assigns to X its picard group, usage \Pic(X)
% \Picardstack_{X/B} denotes the Picard stack of X over B
% \Picardfunctor_{X/B} denotes the Picard functor of X over B
\def\Pic{\mathop{\mathrm{Pic}}\nolimits}
\def\Picardstack{\mathcal{P}\!\mathit{ic}}
\def\Picardfunctor{\mathrm{Pic}}
\def\Deformationcategory{\mathcal{D}\!\mathit{ef}}

\setlength{\emergencystretch}{3em}
\begin{document}
\title[Stacks Project, Tag 0BBJ]{338 --- Blowing up a finite module to make its strict transform locally free}
\maketitle
\noindent Copyright (C) 2005--2025 Johan de Jong. Stacks Project authors. Source: \href{https://stacks.math.columbia.edu/tag/0BBJ}{Tag 0BBJ}.

\noindent Editorial difficulty note (not author text). Difficulty H2: The author reduces a finite module to a finitely presented approximation without losing its strict transform. A Fitting ideal supplies the blowup centre, and nonzerodivisor control of relations proves local freeness of every specified transform; the centre is explicitly constructed..

\noindent Editorial provenance note (not author text). History: excerpt from revision \texttt{a0444\allowbreak 6e57e\allowbreak c1fbc\allowbreak 252a8\allowbreak 71afc\allowbreak ec775\allowbreak 2fb28\allowbreak 07b14}, more-algebra.tex, lines 6228--6297. Reference presentation was adapted; original-author context and core support blocks were retained or added. The mathematical wording is unchanged. Licensed under GNU FDL 1.2 or later; no invariant sections or cover texts. See ../licenses/COPYING. Context blocks reproduce author definitions and supporting results; they are not additional counted problems.

\begin{definition}
\label{definition-strict-transform}
Let $R$ be a ring. Let $I \subset R$ be an ideal and $a \in I$.
Let $R[\frac{I}{a}]$ be the affine blowup algebra, see
Algebra, Definition \ref{algebra-definition-blow-up}.
Let $M$ be an $R$-module.
The {\it strict transform of $M$ along $R \to R[\frac{I}{a}]$} is
the $R[\frac{I}{a}]$-module
$$
M' = \left(M \otimes_R R[\textstyle{\frac{I}{a}}]\right)/a\text{-power torsion}
$$
\end{definition}

\begin{definition}
\label{algebra-definition-blow-up}
Let $R$ be a ring.
Let $I \subset R$ be an ideal.
\begin{enumerate}
\item The {\it blowup algebra}, or the {\it Rees algebra}, associated to
the pair $(R, I)$ is the graded $R$-algebra
$$
\text{Bl}_I(R) =
\bigoplus\nolimits_{n \geq 0} I^n =
R \oplus I \oplus I^2 \oplus \ldots
$$
where the summand $I^n$ is placed in degree $n$.
\item Let $a \in I$ be an element. Denote $a^{(1)}$ the element $a$
seen as an element of degree $1$ in the Rees algebra. Then the
{\it affine blowup algebra} $R[\frac{I}{a}]$ is the algebra
$(\text{Bl}_I(R))_{(a^{(1)})}$ constructed in Section \href{https://stacks.math.columbia.edu/tag/00JM}{section proj (Tag 00JM)}.
\end{enumerate}
\end{definition}

\begin{lemma}
\label{lemma-blowup-module}
Let $R$ be a ring. Let $M$ be a finite $R$-module. Let $f \in R$
be an element such that $M_f$ is finite locally free of rank $r$.
Then there exists a finitely generated ideal $I \subset R$ with
$V(f) = V(I)$ such that for all $a \in I$ with $R' = R[\frac{I}{a}]$
the strict transform
$$
M' = (M \otimes_R R')/a\text{-power torsion}
$$
is locally free of rank $r$.
\end{lemma}

\begin{proof}
Choose a surjection $R^{\oplus n} \to M$. Choose a finite submodule
$K \subset \Ker(R^{\oplus n} \to M)$ such that $R^{\oplus n}/K \to M$
becomes an isomorphism after inverting $f$. This is possible because
$M_f$ is of finite presentation for example by
Algebra, Lemma \href{https://stacks.math.columbia.edu/tag/00NX}{lemma finite projective (Tag 00NX)}.
Set $M_1 = R^{\oplus n}/K$
and suppose we can prove the lemma for $M_1$. Say $I \subset R$ is the
corresponding ideal. Then for $a \in I$ the map
$$
M_1' = (M_1 \otimes_R R')/a\text{-power torsion}
\longrightarrow
M' = (M \otimes_R R')/a\text{-power torsion}
$$
is surjective. It is also an isomorphism after inverting $a$ in $R'$
as $R'_a = R_f$, see Algebra, Lemma \href{https://stacks.math.columbia.edu/tag/080U}{lemma blowup in principal (Tag 080U)}.
But $a$ is a nonzerodivisor on $M'_1$, whence the displayed map is an
isomorphism. Thus it suffices to prove the lemma in case $M$ is a finitely
presented $R$-module.

\medskip\noindent
Assume $M$ is a finitely presented $R$-module.
Then $J = \text{Fit}_r(M) \subset S$ is a finitely generated ideal.
We claim that $I = fJ$ works.

\medskip\noindent
We first check that $V(f) = V(I)$. The inclusion $V(f) \subset V(I)$ is
clear. Conversely, if $f \not \in \mathfrak p$, then
$\mathfrak p$ is not an element of $V(J)$ by 
Lemma \href{https://stacks.math.columbia.edu/tag/07ZC}{lemma fitting ideal generate locally (Tag 07ZC)}.
Thus $\mathfrak p \not \in V(fJ) = V(I)$.

\medskip\noindent
Let $a \in I$ and set $R' = R[\frac{I}{a}]$. We may write $a = fb$
for some $b \in J$. By Algebra, Lemmas \href{https://stacks.math.columbia.edu/tag/07Z3}{lemma affine blowup (Tag 07Z3)} and
\href{https://stacks.math.columbia.edu/tag/0BBI}{lemma blowup add principal (Tag 0BBI)} we see that $J R' = b R'$
and $b$ is a nonzerodivisor in $R'$. 
Let $\mathfrak p' \subset R' = R[\frac{I}{a}]$ be
a prime ideal. Then $JR'_{\mathfrak p'}$ is generated by $b$.
It follows from
Lemma \href{https://stacks.math.columbia.edu/tag/080Z}{lemma principal fitting ideal (Tag 080Z)}
that $M'_{\mathfrak p'}$ can be generated by $r$ elements.
Since $M'$ is finite, there exist $m_1, \ldots, m_r \in M'$ and
$g \in R'$, $g \not \in \mathfrak p'$ such that the corresponding map
$(R')^{\oplus r} \to M'$ becomes surjective after inverting $g$.

\medskip\noindent
Finally, consider the ideal $J' = \text{Fit}_{k - 1}(M')$.
Note that $J' R'_g$ is generated by the coefficients of relations between
$m_1, \ldots, m_r$ (compatibility of Fitting ideal with base change).
Thus it suffices to show that $J' = 0$, see
Lemma \href{https://stacks.math.columbia.edu/tag/07ZD}{lemma fitting ideal finite locally free (Tag 07ZD)}.
Since $R'_a = R_f$ (Algebra, Lemma \href{https://stacks.math.columbia.edu/tag/080U}{lemma blowup in principal (Tag 080U)})
and $M'_a = M_f$ is free of rank $r$ we see that $J'_a = 0$.
Since $a$ is a nonzerodivisor in $R'$ we
conclude that $J' = 0$ and we win.
\end{proof}
\end{document}
